% Source: https://artofproblemsolving.com/community/c3h1580925_geometry
% Stable identity: aops-topic-1580925-q01
Since $3^2+4^2=5^2$, the base is a right triangle with legs $3$ and $4$, so its area is $\dfrac{1}{2}\cdot3\cdot4=6$ square feet. The volume of a pyramid is one third of the base area times the height:
\[
V=\dfrac{1}{3}\cdot6\cdot6=\boxed{12}\text{ cubic feet}.
\]

\begin{asy}
usepackage("amsmath");
size(5cm);
defaultpen(fontsize(10pt));
pair P3(real x,real y,real z){return (x+0.45*y,z+0.3*y);}
pair R=P3(0,0,0),Pp=P3(4,0,0),Q=P3(0,3,0),T=P3(4/3,1,6),F=P3(4/3,1,0);
draw(R--Pp^^R--T^^Pp--T,linewidth(0.9));
draw(R--Q^^Q--Pp^^Q--T,dashed);
draw(T--F,red+dashed);
real s=0.3; draw(P3(0,s,0)--P3(s,s,0)--P3(s,0,0));
label("$4$",(R+Pp)/2,S);label("$3$",(R+Q)/2,NW);label("$5$",(Q+Pp)/2,NE);
label("$6$",(T+F)/2,E,red);
dot(T);dot(F,red);
label("\mbox{}",(min(currentpicture,user=true)+max(currentpicture,user=true))/2,black); // keeps the text after the figure black when LaTeX typesets the labels inline
\end{asy}

(The sketch is an oblique view; hidden edges are dashed and the red dashed segment is the altitude. Where the foot of the altitude lies does not matter.)
